\subsection{Traffic Prediction}
\label{sec:traffic-prediction}

GetS requests and DataS responses use XY routing on plane C. The host
optimizes MoE routes on plane M. For Dispatch, we adapt MixNet's
adjacent-layer predictor~\cite{liao2025mixnet}:
\begin{equation}
  \widehat{\mathbf y}_{\ell,s}=P_\ell\mathbf x_{\ell-1,s},
  \quad P_\ell\geq0,\quad
  \mathbf1^{\mathsf T}P_\ell=\mathbf1^{\mathsf T}.
  \label{eq:noc-load-prediction}
\end{equation}
Here $\mathbf x_{\ell-1,s}$ is the preceding layer's normalized expert distribution
for tokens with known source NI $s$. Weighted least squares fits $P_\ell$
from paired token groups, weighting their size and recency. Prefill and
decode use separate models, with a historical-mean fallback.
Expert placement and packing map forecasts to source--destination
demand. Final allocation uses actual Top-K and packing metadata.
Combine uses confirmed assignments, output shapes, and estimated
expert completion times to derive return packets and release times.

\subsection{Reconfiguration Problem}
\label{sec:noc-problem}

Within a fixed $H$-cycle planning interval, $\mathcal M_f$ contains
the unbound packets of MoE source--destination flow $f$.
Their payloads $b_m$ give $D_f=\sum_{m\in\mathcal M_f}b_m$ and
$\mathcal M=\bigcup_f\mathcal M_f$. Legal CS paths for packet $m$
form $\mathcal P_m$, with $\mathcal P=\bigcup_m\mathcal P_m$.
Each simple path connects the source and destination NI through wired
CS ports. Binary $z_p$ selects path $p$;
binary $x_{m,p}$ assigns packet $m$ to it. We maximize useful payload
assigned to the CS network:
\begin{equation}
  \max_{\mathbf z,\mathbf x}\quad
  V=\sum_{m\in\mathcal M}b_m\sum_{p\in\mathcal P_m}x_{m,p}.
  \label{eq:noc-volume-objective}
\end{equation}
Readiness flag $g_m$ equals one when source data and receive storage
are ready, and zero otherwise. Each packet uses at most one enabled path:
\begin{equation}
  \sum_{p\in\mathcal P_m}x_{m,p}\leq1,
  \qquad x_{m,p}\leq z_p,\quad x_{m,p}\leq g_m.
  \label{eq:noc-packet-assignment}
\end{equation}
Selected paths are CS-disjoint. With $A_{c,p}=1$ when $p$ uses CS $c$,
\begin{equation}
  \sum_{p\in\mathcal P}A_{c,p}z_p\leq1,
  \qquad c=0,\ldots,12.
  \label{eq:noc-cs-exclusion}
\end{equation}

Resource $r\in\mathcal R$ is an M-plane link or NI, DMA, or SRAM service.
Its remaining budget $C_r$ within $H$ deducts work outside $\mathcal M$:
existing transfers, configuration messages, and C-plane use of shared endpoints.
Cost $u_{m,p,r}$ includes framing and expresses service in NoC cycles.
Unassigned packets use the M-plane Mesh, denoted
by 0, with $x_{m,0}=1-\sum_{p\in\mathcal P_m}x_{m,p}$. Thus,
\begin{equation}
  \sum_{m\in\mathcal M}\sum_{p\in\mathcal P_m\cup\{0\}}
  u_{m,p,r}x_{m,p}\leq C_r,\qquad r\in\mathcal R.
  \label{eq:noc-resource-capacity}
\end{equation}
At most two paths are selected, $\sum_{p\in\mathcal P}z_p\leq2$;
a path may traverse multiple CSs.
Path $p$ specifies the input, output, and enable fields of each CS;
$x_{m,p}$ determines the source NI's route tag for packet $m$.

\subsection{Configuration Algorithm}
\label{sec:noc-algorithm}

Offline enumeration supplies local and upper-CS paths. We target the two
M-plane links with the largest pending MoE flit counts. For each link
still unrelieved, flows are visited in decreasing $D_f$, with endpoint
identifiers breaking ties. Up to $K_p=2$ paths per flow are tried by length
and then identifier. We retain the first CS-disjoint extension that passes
joint \textsc{Assign} and increases $V$. Construction ends after two paths,
both links, or budget exhaustion. At most two configurations are retained;
the one-path configuration remains available when a second path is added.

Candidates are checked in decreasing $V$. A timing model propagates
source reads, finite queues, credits, receive writes, and dependencies;
DataS follows its request's memory service. While $S_0$ forwards, an update
must satisfy, using the same time origin,
\begin{equation}
  \widehat T(S)+\eta(S)<\widehat T(S_0).
  \label{eq:noc-switching-check}
\end{equation}
Here $\eta(S)\geq0$ is calibrated from paired errors on held-out requests.
Protected consumers satisfy
$\widehat T_k(S)-\widehat T_k(S_0)+\eta_k(S)\leq\epsilon_k$,
including MoE with $\epsilon_k=0$ and coherence transactions with their
specified allowances. Estimates include controller execution, pause,
drain, writes, and NI installation. After pause acknowledgement and drain, versions and eligibility are refreshed and feasibility, timing, and budgets are
rechecked before Section~\ref{sec:noc-config-interface}'s update protocol.

\begin{algorithm}[H]
  \caption{CS configuration for at most two hotspots}
  \label{alg:noc-configuration}
  \small
  \begin{algorithmic}[1]
    \Require MoE packets $\mathcal M$, paths $\mathcal P$, budgets $C_r$, active $S_0$
    \Ensure $S_0$, an acknowledged update, or an update error
    \State $\mathcal H\gets$ up to two M-plane links with largest demand
    \State $S\gets$ Mesh allocation; $\mathcal C\gets\varnothing$
    \For{$e\in\mathcal H$ while construction budget remains}
      \State \textbf{if} $S$ already relieves $e$ \textbf{then continue}
      \State $S'\gets$ first feasible extension for $e$ increasing $V$, or $S$
      \If{$V(S')>V(S)$}
        \State $S\gets S'$; $\mathcal C\gets\mathcal C\cup\{S\}$
      \EndIf
      \State \textbf{if} two paths are selected \textbf{then break}
    \EndFor
    \For{$S'\in\mathcal C$ in decreasing $V$}
      \State \textbf{if} budget expired \textbf{then return} $S_0$
      \State Estimate completion and update overhead
      \If{timing checks pass within the decision budget}
        \State Request packet holds and pause new CS admissions
        \If{pause and drain complete within budget}
          \State Refresh state; recheck feasibility and timing
          \State \textbf{if} final checks pass \textbf{then return} $\Call{ApplyCS}{S'}$
        \EndIf
        \State Release holds and resume $S_0$; \Return $S_0$
      \EndIf
    \EndFor
    \State \Return $S_0$
  \end{algorithmic}
\end{algorithm}
